\begin{proof}
Para demostrar la equivalencia de \emph{a)}, \emph{b)} y \emph{c)}, basta evidentemente establecer
que \emph{b)} implica \emph{c)}. Sea, pues, $K$ una extensión de tipo finito de $k$, que es
por tanto el cuerpo de fracciones de una $k$-álgebra de tipo finito $A$; pongamos
$Y=\operatorname{Spec}(A)$. Se sabe (4.6.6) que existe una extensión finita
y radicial $k'$ de $k$ tal que $Y'=(Y\otimes_k k')_{\mathrm{red}}$ sea
separable sobre $k'$; si $\eta'$ es el punto genérico de $Y'$, $K'=\boldsymbol{k}(\eta')$
es una extensión separable de $k'$ por (4.6.1). Dicho esto,
$\boldsymbol{P}(X_{(k')},x')$ es verdadera por hipótesis para todo punto $x'$ de $X_{(k')}$
situado sobre $x$, por lo que resulta de (6.7.4) que $\boldsymbol{P}(X_{(K')},x'')$ es
verdadera para todo punto $x''$ de $X_{(K')}$ situado sobre $x$, puesto que
$X_{(K')}=(X_{(k')})_{(K')}$, que $x''$ está situado sobre un punto $x'$ de
$X_{(k')}$, él mismo situado sobre $x$, y que $K'$ es separable sobre $k'$.
Pero también se tiene
$X_{(K')}=(X_{(K)})_{(K')}$ y para todo $x_0\in X_{(K)}$ situado sobre
$x$, existe $x''\in X_{(K')}$ situado sobre $x_0$; resulta, pues, de
(6.7.4) que $\boldsymbol{P}(X_{(K)},x_0)$ es verdadera.

Supongamos ahora que $X$ es localmente de tipo finito sobre $k$, de modo que
para \emph{toda} extensión $K$ de $k$, $X_{(K)}$ es localmente noetheriano. Como una
extensión radicial de $k$ es isomorfa a una subextensión de cualquier extensión
perfecta de $k$, resulta inmediatamente de (6.7.4) que \emph{e)} implica \emph{f)}; del mismo
modo, toda extensión finita radicial de $k$ está contenida en una extensión
$k^{p^{-s}}\subset k^{p^{-\infty}}$, por lo que \emph{f)} implica \emph{b)}; \emph{d)} implica
trivialmente \emph{e)}, y finalmente \emph{e)} implica \emph{d)}. En efecto, si $K$ es una extensión
\emph{cualquiera} de $k$,
existe una extensión $K'$ de $k$ que contiene (salvo $k$-isomorfismos) a
$k'$ y a $K$; como $K'$ es una extensión \emph{separable} de $k'$, se deduce de (6.7.4)
que $\boldsymbol{Q}(k')$ es verdadera, y después que $\boldsymbol{Q}(K)$ es verdadera, razonando como en la
primera parte de la demostración. Basta, pues, demostrar que \emph{b) implica
e)}. Tomaremos $k'=k^{p^{-\infty}}$ en \emph{e)}.

Como la cuestión es local sobre $X$, se puede además restringir al caso en que $X$
es afín y de tipo finito sobre $k$, reemplazando $X$ por una vecindad de $x$;
sea, pues, $X=\operatorname{Spec}(B)$, donde $B$ es una $k$-álgebra de tipo finito;
además, por definición de la propiedad $\boldsymbol{P}$ en los casos considerados, también se
puede reemplazar $X$ por el esquema local de $X$ en el punto $x$, y suponer, por
tanto, que $B$ es un anillo \emph{local} noetheriano. Pongamos $B'=B\otimes_k k'$; $k'$
es límite inductivo de sus subextensiones \emph{finitas} $k'_\lambda$ y si se pone
$B'_\lambda=B\otimes_k k'_\lambda$, el morfismo
$f_\lambda:\operatorname{Spec}(B')\to\operatorname{Spec}(B'_\lambda)$ es
un \emph{homeomorfismo} de $\operatorname{Spec}(B')$ \emph{sobre}
$\operatorname{Spec}(B'_\lambda)$ para todo $\lambda$;
\oldpage[IV-2]{150}
en efecto, $f_\lambda$ es biyectivo en virtud de (\textbf I, 3.5.2, 3.5.7 y 3.5.8);
por otra parte, es cerrado por (\textbf{II}, 6.1.10). Ahora se puede aplicar
(5.13.6), que demuestra (en virtud de la hipótesis \emph{b)}) que \emph{b)} implica \emph{e)} cuando
$\boldsymbol{P}(Z,z)$ es la propiedad $(R_n)$ en el punto $z$. El caso en que $\boldsymbol{P}(Z,z)$ es la
propiedad de ser regular (es decir, de poseer la propiedad $(R_n)$ para todo $n$)
resulta trivialmente de ello. Finalmente, teniendo en cuenta el criterio de
normalidad de Serre (5.8.6), \emph{b)} implica también \emph{e)} cuando $\boldsymbol{P}(Z,z)$ es la propiedad
de ser normal en el punto $z$, pues esto equivale a decir que $Z$ posee en el
punto $z$ las propiedades $(S_2)$ y $(R_1)$: se aplica entonces lo que precede
para $n=1$, y (6.7.2, (ii)) para $n=2$.
\end{proof}
